\scanpage{39}
Donc $U_\rho$ opère \inserted{\uline{fidèlement}} sur $\mathcal E,\mathcal E'$ par
\[
\left\{
\begin{aligned}
u_\lambda^{\mathcal E}&=\operatorname{id}+\lambda\Pi'\otimes\Pi\\
u_\lambda^{\mathcal E'}&=\operatorname{id}+\lambda\Pi'\otimes\Pi
\end{aligned}\right.
\]
(opérations \deleted{\unsure{contragrédientes}}\inserted{transposées} l'une de l'autre)
Comme $\Pi,\Pi'$ invariants sous les $\tau_i$, donc sous les
$\varphi^{\mathcal E}(s)$, $\varphi_{\unsure{s}}^{\mathcal E'}(s)$
$(s\in\unsure{\mathfrak G_n})$, on trouve
\[
\unsure{(\alpha')}\qquad
\diagram{\begin{tikzcd}[ampersand replacement=\&]
U_\rho\arrow[r] \& \uline{\operatorname{Aut}}(\mathcal E,\varphi^{\mathcal E})
\end{tikzcd}}
\]
\[
\unsure{(\alpha)}\qquad
\diagram{\begin{tikzcd}[ampersand replacement=\&]
U_\rho\arrow[r] \& \uline{\operatorname{Aut}}(X,\varphi^X)
\end{tikzcd}}
\]
et de $=$ sur $\mathcal E'$. On obtient en fait, en \unsure{adjoignant} les scalaires
\[
\unsure{(\alpha\alpha)}\qquad
\diagram{\begin{tikzcd}[ampersand replacement=\&,column sep=large]
\mathbb G_m\times U_\rho\arrow[rr]\arrow[dr,bend right]\& \&
\operatorname{Aut}(\mathcal E,\varphi^{\mathcal E})\\
\&\mathbb G_m\arrow[ur,"\text{opère par \unsure{homothéties}}"']\&
\end{tikzcd}}
\]
dont $\unsure{\alpha}$ est déduit en passant au quotient. On trouve

\uline{Thm}\quad $\unsure{\alpha}$ et $\unsure{(\alpha\alpha)}$ sont des iso.

\margin{\figure{\begin{tikzpicture}\draw (0,0)--(0,-1.7);\end{tikzpicture}}}
De plus, si $D_0=(H_1\cap\cdots\cap H_n)$,
$\uline{D_0^{**}}=D_0\setminus\{c,d\}$, où
$\unsure{\{d\}}=D_0\cap C$ (qui est une section de $X$ sur $S$) on a que
\mspage{38}

